\documentclass[12pt,a4paper]{article}

\usepackage[margin=2.5cm]{geometry}
\usepackage{fontspec}
\setmainfont{Latin Modern Roman}
\setsansfont{Latin Modern Sans}
\setmonofont{Latin Modern Mono}
\usepackage{amsmath,amssymb,amsthm,mathtools,bm}
\usepackage{microtype}
\usepackage{booktabs}
\usepackage{array}
\usepackage{enumitem}
\usepackage{hyperref}
\usepackage{xcolor}
\usepackage{fancyhdr}
\usepackage{lastpage}

\hypersetup{
  colorlinks=true,
  linkcolor=black,
  citecolor=black,
  urlcolor=blue,
  pdftitle={The Space of Geometry: From First Distinction to Pythagoras},
  pdfauthor={Adrian Lipa},
  pdfsubject={Theory of Informational Relations; affine quantum-state geometry; Euclidean closure; tetrahedral simplex},
  pdfkeywords={TIR, geometry, qubit, affine space, Hermitian operators, tetrahedron, Pythagoras, SIC-POVM}
}

\pagestyle{fancy}
\fancyhf{}
\lhead{The Space of Geometry}
\rhead{Adrian Lipa}
\cfoot{\thepage\ / \pageref{LastPage}}
\setlength{\headheight}{15pt}

\newtheorem{theorem}{Theorem}[section]
\newtheorem{lemma}[theorem]{Lemma}
\newtheorem{proposition}[theorem]{Proposition}
\newtheorem{corollary}[theorem]{Corollary}
\newtheorem{definition}[theorem]{Definition}
\newtheorem{remark}[theorem]{Remark}

\newcommand{\Herm}{\operatorname{Herm}}
\newcommand{\Aff}{\operatorname{Aff}}
\newcommand{\Tr}{\operatorname{Tr}}
\newcommand{\Aut}{\operatorname{Aut}}
\newcommand{\R}{\mathbb{R}}
\newcommand{\C}{\mathbb{C}}
\newcommand{\I}{\mathbb{I}}
\newcommand{\Erel}{\mathcal{E}}
\newcommand{\Astate}{\mathcal{A}_2}
\newcommand{\sigmavec}{\bm{\sigma}}
\newcommand{\rvec}{\bm{r}}
\newcommand{\nvec}{\bm{n}}
\newcommand{\vvec}{\bm{v}}

\title{\textbf{The Space of Geometry}\\[0.4em]\large From First Distinction to Pythagoras}
\author{Adrian Lipa\\\small CIEL/0 Project -- Intention Lab}
\date{August 2026}

\begin{document}
\maketitle

\begin{abstract}
We construct a local relation geometry from the binary quantum-point carrier of the Theory of Informational Relations (TIR). The density-operator set of a two-level quantum system has affine hull
\[
\Astate=\frac12 I+\Herm_0(2),
\]
whose translation space has real dimension three. Every ordered pair of normalized point states determines the unique affine displacement
\[
\delta(\rho_x,\rho_y)=\rho_y-\rho_x.
\]
With the Pauli normalization used in the TIR spatial branch,
\[
\Erel_{xy}=2(\rho_y-\rho_x)\in\Herm_0(2)\cong\R^3.
\]
The conjugation action of $SU(2)$ induces the standard $SO(3)$ action on this real carrier and preserves the inner product $\langle A,B\rangle=\Tr(AB)/2$. The construction then forks. The Euclidean branch combines additive endpoint composition with the invariant inner product and reaches the Pythagorean identity for orthogonal consecutive relations. The finite-cell branch asks for the smallest full-dimensional affine support of the same three-dimensional carrier: four affinely independent vertices give a $3$-simplex. Its unlabeled automorphism group is $S_4$, transitive on the six edges; applying the TIR arithmetic-geometric measure axiom A5 and law-symmetry axiom A7 to the edge-measure law yields equal edge lengths and hence the regular tetrahedron. Its centered unit frame obeys $\nvec_a\cdot\nvec_b=-1/3$ for $a\neq b$. Independently, the minimal symmetric qubit informationally complete rank-one measurement has the same tetrahedral Gram geometry. Thus the local Pythagorean closure and the tetrahedral finite-cell theorem are parallel consequences of one canonical three-real-dimensional relation carrier.
\end{abstract}

\section{Question and scope}

The question is deliberately narrow:
\begin{center}
\fbox{\parbox{0.86\linewidth}{\centering\small How little primitive relational structure is required for local Euclidean geometry to appear?}}
\end{center}
The principal endpoint is
\[
\boxed{a^2+b^2=c^2.}
\]
The paper also derives the minimal finite full-dimensional local cell and shows how its regularity follows from the TIR measurement and symmetry axioms. The two results share a common carrier but form separate theorem branches.

The dependency structure is
\[
\boxed{
\C^2
\longrightarrow
\rho_x
\longrightarrow
\Astate
\longrightarrow
\delta(\rho_x,\rho_y)
\longrightarrow
V=\Herm_0(2)\cong\R^3
}
\]
followed by
\[
\boxed{
V\longrightarrow
\begin{cases}
\text{invariant inner-product geometry}\longrightarrow\text{Pythagoras},\\[2mm]
\text{minimal finite support}\longrightarrow\Delta^3\xrightarrow{A5+A7}\text{regular tetrahedron}.
\end{cases}}
\]
A third, information-theoretic branch reaches the same tetrahedral Gram frame through symmetric informational completeness.

\section{Imported TIR root}

The parent TIR programme supplies the primitive relational chain
\[
0\longrightarrow P\longrightarrow\text{first distinction}\longrightarrow\{N,S\}
\longrightarrow\frac12\longrightarrow\ln 2\longrightarrow\C^2.
\]
The present derivation begins its new geometric work at the binary quantum carrier. The first distinction provides two quantum alternatives, and their coherent span is the minimal two-level Hilbert carrier $\mathcal H_2\cong\C^2$.

A normalized two-level quantum state is represented by a density operator
\[
\rho_x=\rho_x^\dagger,\qquad \Tr\rho_x=1,\qquad \rho_x\ge 0.
\]
Every such state has Bloch form \cite{NielsenChuang}
\[
\boxed{
\rho_x=\frac12\left(I+\rvec_x\cdot\sigmavec\right),\qquad |\rvec_x|\le 1,
}
\]
where $\sigmavec=(\sigma_x,\sigma_y,\sigma_z)$ is the Pauli triple. The physical density operators fill the Bloch ball, while their affine hull is the full trace-one Hermitian hyperplane used below.

\section{The common three-real-dimensional carrier}

\begin{theorem}[Affine carrier theorem]\label{thm:carrier}
The affine hull of normalized Hermitian $2\times2$ states is
\[
\Astate=\frac12 I+\Herm_0(2),
\]
and its translation space has real dimension three:
\[
V=\Herm_0(2)=\operatorname{span}_{\R}\{\sigma_x,\sigma_y,\sigma_z\}\cong\R^3.
\]
\end{theorem}

\begin{proof}
Every Hermitian $2\times2$ matrix can be written uniquely as
\[
A=a_0I+a_x\sigma_x+a_y\sigma_y+a_z\sigma_z,
\qquad a_\mu\in\R.
\]
Thus $\dim_{\R}\Herm(2)=4$. Since $\Tr I=2$ and $\Tr\sigma_i=0$, the affine trace-one condition fixes $a_0=1/2$. Differences of any two trace-one Hermitian operators are traceless Hermitian, hence the translation space is $\Herm_0(2)$. The Pauli matrices form a real basis of that space, so
\[
\boxed{\dim_{\R}\Herm_0(2)=3.}
\]
\end{proof}

The number three enters this branch as the real dimension of the translation carrier of the normalized binary quantum-state hull.

\section{Canonical endpoint-carrying relation}

An affine space is a torsor for its translation space: for every ordered pair $x,y$ there is a unique translation vector carrying $x$ to $y$.

\begin{definition}[Canonical affine displacement]
For $\rho_x,\rho_y\in\Astate$, define
\[
\boxed{\delta_{xy}:=\delta(\rho_x,\rho_y)=\rho_y-\rho_x.}
\]
The TIR spatial branch exports this endpoint-carrying torsor displacement as its primitive vector relation. In Pauli generator normalization,
\[
\boxed{
\Erel_{xy}:=2\delta_{xy}=2(\rho_y-\rho_x)
=(\rvec_y-\rvec_x)\cdot\sigmavec.
}
\]
\end{definition}

\begin{proposition}[Endpoint identities]\label{prop:endpoint}
The canonical relation satisfies
\[
\Erel_{yx}=-\Erel_{xy},
\]
\[
\boxed{\Erel_{xz}=\Erel_{xy}+\Erel_{yz},}
\]
and
\[
\boxed{\Erel_{xy}+\Erel_{yz}+\Erel_{zx}=0.}
\]
\end{proposition}

\begin{proof}
All three identities are immediate from telescoping differences of the endpoints. For example,
\[
2(\rho_y-\rho_x)+2(\rho_z-\rho_y)=2(\rho_z-\rho_x).
\]
\end{proof}

The construction uses only the ordered endpoints and the affine structure already present in the state carrier. The earlier TIR source-minimality/naturality result gives an independent uniqueness crosscheck: any continuous affine-natural vector relation built from the ordered endpoints alone has the form $c(y-x)$, with the nonzero scalar fixing the global relation unit and orientation.

\section{Rotational covariance and the invariant metric}

For $U\in SU(2)$, a common quantum-frame transformation acts by conjugation,
\[
\rho\mapsto U\rho U^\dagger.
\]
Therefore
\[
\boxed{\Erel_{xy}\mapsto U\Erel_{xy}U^\dagger.}
\]
The induced action on the three real Pauli coefficients factors through the standard double cover
\[
\boxed{SU(2)/\{\pm I\}\cong SO(3).}
\]

The natural positive invariant bilinear form on $\Herm_0(2)$ is
\[
\boxed{\langle A,B\rangle:=\frac12\Tr(AB).}
\]
It obeys
\[
\frac12\Tr(\sigma_i\sigma_j)=\delta_{ij},
\]
so the Pauli basis is orthonormal. Conjugation preserves the form because of cyclicity of trace:
\[
\Tr\left((UAU^\dagger)(UBU^\dagger)\right)=\Tr(AB).
\]
Since the defining real $SO(3)$ representation is irreducible, an invariant positive inner product is unique up to an overall positive scale. The above convention fixes that scale.

For the canonical relation,
\[
\boxed{
\|\Erel_{xy}\|^2
=\frac12\Tr(\Erel_{xy}^2)
=|\rvec_y-\rvec_x|^2.
}
\]
This is the arithmetic-geometric quantity used by A5 as the squared local relation length.

\section{Theorem E: local Euclidean and Pythagorean closure}

\begin{theorem}[Theorem E]\label{thm:pyth}
Let
\[
A=\Erel_{xy},\qquad B=\Erel_{yz},
\]
be two consecutive local relations in one affine comparison frame. Then endpoint composition gives $\Erel_{xz}=A+B$. With the invariant inner product $\langle A,B\rangle=\Tr(AB)/2$, the carrier supports the Euclidean norm identity
\[
\|A+B\|^2=\|A\|^2+\|B\|^2+2\langle A,B\rangle.
\]
For orthogonal relations, $\langle A,B\rangle=0$, and therefore
\[
\boxed{a^2+b^2=c^2},
\]
where $a=\|A\|$, $b=\|B\|$, and $c=\|A+B\|$.
\end{theorem}

\begin{proof}
By Proposition~\ref{prop:endpoint}, $\Erel_{xz}=A+B$. Bilinearity and symmetry of the inner product give
\begin{align*}
\|A+B\|^2
&=\langle A+B,A+B\rangle\\
&=\langle A,A\rangle+\langle B,B\rangle+2\langle A,B\rangle.
\end{align*}
For $A\perp B$, the cross term vanishes. Setting $a=\|A\|$, $b=\|B\|$, and $c=\|A+B\|$ yields the Pythagorean identity.
\end{proof}

For nonzero relations, the local angle is consequently defined by
\[
\boxed{
\cos\theta=\frac{\langle A,B\rangle}{\|A\|\,\|B\|},
}
\]
and orthogonality by
\[
A\perp B\iff\langle A,B\rangle=0.
\]
Theorem E is the principal endpoint of the paper.

\section{Theorem T1: minimal finite full-dimensional support}

The same common carrier also answers a different question: what is the smallest finite affine cell that can support the derived three-dimensional relation space?

\begin{theorem}[Theorem T1]\label{thm:simplex}
A finite full-dimensional affine cell in a three-dimensional carrier requires at least four affinely independent vertices. At the minimum it is a $3$-simplex,
\[
\boxed{\Delta^3.}
\]
\end{theorem}

\begin{proof}
For points $x_1,\ldots,x_m$ in an affine space,
\[
\dim\Aff\{x_1,\ldots,x_m\}\le m-1.
\]
Full support of a three-dimensional affine carrier requires
\[
3\le m-1,
\]
so $m\ge4$. At equality, four affinely independent points define precisely the simplex
\[
\Delta^3=\operatorname{conv}\{x_1,x_2,x_3,x_4\}.
\]
\end{proof}

Thus the tetrahedron appears as the minimal finite full-dimensional affine cell of the same local carrier that already supports Theorem E.

\section{Theorem T2: regularity from A5+A7 edge-orbit invariance}

The abstract unlabeled $3$-simplex has automorphism group
\[
\boxed{\Aut(\Delta^3)\cong S_4.}
\]
Its six edges are the unordered vertex pairs $\{i,j\}$ with $1\le i<j\le4$. The natural $S_4$ action on these six pairs is transitive; therefore the edges form one intrinsic combinatorial orbit.

For every realized edge define the A5 arithmetic-geometric measure
\[
\boxed{
q_{ij}:=\frac12\Tr(\Erel_{ij}^2),\qquad
\ell_{ij}:=\sqrt{q_{ij}}.
}
\]

\begin{theorem}[Theorem T2]\label{thm:regular}
Apply A7 law symmetry to the intrinsic unlabeled-simplex automorphism action, so that the edge-measure law satisfies
\[
q_{\pi(i)\pi(j)}=q_{ij}
\qquad\text{for every }\pi\in S_4.
\]
Then every edge has the same measure and the same length. The minimal $3$-simplex is therefore regular.
\end{theorem}

\begin{proof}
Transitivity means that for any two edges $e$ and $e'$ there exists $\pi\in S_4$ with $\pi(e)=e'$. A7 invariance then gives $q_{e'}=q_e$. Hence there is one common value $q_*$ on all six edges, and consequently one common length $\ell_*$. A nondegenerate Euclidean tetrahedron with six equal edges is regular.
\end{proof}

The minimal-cell branch is therefore
\[
\boxed{
\Herm_0(2)\cong\R^3
\longrightarrow
\Delta^3
\xrightarrow{A5+A7}
\text{regular tetrahedron}.
}
\]

\section{Tetrahedral Gram geometry and the appearance of \texorpdfstring{$-1/3$}{-1/3}}

Place the barycenter of a regular tetrahedron at the origin and normalize the four center-to-vertex directions $\nvec_a$ so that $|\nvec_a|=1$. The barycenter condition gives
\[
\boxed{\sum_{a=1}^{4}\nvec_a=0.}
\]
Regularity makes every off-diagonal inner product equal to a common value $q$. Hence
\[
0=\left|\sum_{a=1}^{4}\nvec_a\right|^2
=4+12q,
\]
which yields
\[
\boxed{q=-\frac13.}
\]
Therefore
\[
\boxed{\nvec_a\cdot\nvec_b=-\frac13\qquad(a\neq b).}
\]
The Gram matrix is
\[
\boxed{
G=\frac43 I_4-\frac13\mathbf 1\mathbf 1^{T},
}
\]
and the frame operator is
\[
\boxed{
\sum_{a=1}^{4}\nvec_a\nvec_a^{T}=\frac43 I_3.
}
\]
The orientation-preserving rotational symmetry group is $A_4\subset SO(3)$, whose inverse image under the $SU(2)$ double cover is the binary tetrahedral group $2T\subset SU(2)$ \cite{Coxeter}.

This separates two fractions that arise at different structural levels of the construction. The value $1/2$ is the normalized fixed share of the primitive binary distinction in the parent TIR root. The value $-1/3$ is the off-diagonal Gram invariant of the regular tetrahedral frame in the derived three-real-dimensional carrier:
\[
\boxed{
\frac12\quad\longrightarrow\quad\C^2\quad\longrightarrow\quad\R^3\quad\longrightarrow\quad-\frac13.
}
\]

\section{Theorem Q: independent tetrahedral informational convergence}

A general qubit density operator carries three independent real Bloch coordinates. A normalized probability vector with $m$ outcomes carries at most $m-1$ independent real entries. Hence informational completeness requires
\[
m-1\ge3,
\]
so
\[
\boxed{m\ge4.}
\]
For the symmetric rank-one minimal case, the qubit SIC-POVM is tetrahedral \cite{RenesEtAl}. Let $\nvec_a$ be the same type of regular tetrahedral unit frame, and define
\[
\boxed{
F_a=\frac14\left(I+\nvec_a\cdot\sigmavec\right),\qquad a=1,\ldots,4.
}
\]
Because $\sum_a\nvec_a=0$, the effects sum to the identity:
\[
\sum_{a=1}^{4}F_a=I.
\]
For
\[
\rho=\frac12\left(I+\rvec\cdot\sigmavec\right),
\]
the outcome probabilities are
\[
\boxed{
p_a=\Tr(\rho F_a)=\frac14\left(1+\rvec\cdot\nvec_a\right).
}
\]
Using the tight-frame identity $\sum_a\nvec_a\nvec_a^T=(4/3)I_3$,
\begin{align*}
\sum_a p_a\nvec_a
&=\frac14\sum_a\nvec_a+\frac14\sum_a(\rvec\cdot\nvec_a)\nvec_a\\
&=\frac14\left(\frac43 I_3\right)\rvec
=\frac13\rvec,
\end{align*}
so the state reconstructs as
\[
\boxed{\rvec=3\sum_{a=1}^{4}p_a\nvec_a.}
\]

\begin{theorem}[Theorem Q]
The minimal symmetric rank-one informationally complete measurement for a qubit has four outcomes and a tetrahedral Bloch frame with pairwise Gram value $-1/3$. It therefore converges on the same finite frame geometry as the regular minimal spatial simplex derived in Theorems~\ref{thm:simplex}--\ref{thm:regular}.
\end{theorem}

The spatial edge relation and the informational measurement outcome retain distinct physical typing; the exact equality here is the finite Gram-frame convergence.

\section{Proof-status and dependency audit}

The role of each step is summarized in Table~\ref{tab:status}. This keeps standard mathematics, TIR axiomatic input, and independent convergence results visibly separated.

\begin{table}[htbp]
\centering
\small
\renewcommand{\arraystretch}{1.25}
\begin{tabular}{>{\raggedright\arraybackslash}p{0.43\textwidth} >{\raggedright\arraybackslash}p{0.47\textwidth}}
\toprule
\textbf{Result} & \textbf{Status / dependency}\\
\midrule
Binary quantum carrier $\C^2$ & Imported TIR root from the first distinction and A2 quantum-point postulate\\
$\Astate=I/2+\Herm_0(2)$ & Exact linear algebra / standard two-level quantum representation\\
$\dim_{\R}\Herm_0(2)=3$ & Exact linear algebra\\
$\delta_{xy}=\rho_y-\rho_x$ & Exact affine-torsor displacement; exported as the TIR spatial vector relation\\
Endpoint composition and triangular closure & Exact affine identities\\
$SU(2)/\{\pm I\}\cong SO(3)$ and trace metric & Exact representation theory / matrix algebra\\
Theorem E: Pythagorean closure & Exact inner-product consequence on the local relation carrier\\
Theorem T1: $m\ge4$ and $\Delta^3$ & Exact affine-dimension theorem\\
Theorem T2: regular tetrahedron & Exact conditional consequence of A5 measure plus A7 invariance on the intrinsic $S_4$ edge orbit\\
Tetrahedral Gram value $-1/3$ & Exact consequence of regularity and barycentric normalization\\
Theorem Q: tetrahedral qubit SIC convergence & Standard quantum-information construction, independently convergent with the spatial finite frame\\
Global triangulation, curvature, holonomy, torsion, scale calibration, TIR $\times$ Time closure & Downstream research programme\\
\bottomrule
\end{tabular}
\caption{Dependency and proof-status surface for the local construction.}
\label{tab:status}
\end{table}

\section{Global geometry boundary}

For a regular tetrahedron, the dihedral angle is
\[
\boxed{\theta_T=\arccos\left(\frac13\right).}
\]
It satisfies
\[
\boxed{5\theta_T<2\pi<6\theta_T.}
\]
Thus the local regular tetrahedral frame naturally opens a separate global refinement problem. Piecewise-flat tetrahedral complexes can encode curvature through angular deficits, while transport and nontrivial holonomy enter at the gluing level. Physical scale calibration and the TIR $\times$ Time spacetime join belong to the same downstream programme.

\section{Conclusion}

The local construction begins with a binary quantum carrier and uses the affine geometry already present in its normalized state hull. Trace normalization leaves the translation space
\[
\Herm_0(2)\cong\R^3,
\]
and the unique endpoint-carrying torsor displacement provides an origin-independent relation vector. The $SU(2)$ conjugation action induces the standard $SO(3)$ rotational action, while the trace bilinear form supplies the invariant Euclidean inner product.

From this common carrier, two parallel theorem branches follow. The first is immediate: additive endpoint composition plus orthogonality yields
\[
\boxed{a^2+b^2=c^2.}
\]
The second asks for minimal finite full-dimensional support. Three real affine dimensions require four vertices, giving $\Delta^3$; A5 measurement and A7 symmetry on the single intrinsic $S_4$ edge orbit select equal edge lengths and therefore the regular tetrahedron. Its centered unit frame has the characteristic Gram value $-1/3$. A minimal symmetric qubit SIC reaches the same tetrahedral frame independently.

The final local structure is
\[
\boxed{
\begin{aligned}
\text{binary quantum point}
&\longrightarrow \text{canonical 3D relation carrier}\\
&\longrightarrow
\begin{cases}
\text{Euclidean branch}\longrightarrow\text{Pythagoras},\\
\text{finite cell}\longrightarrow\text{regular tetrahedron}\longleftarrow\text{SIC}.
\end{cases}
\end{aligned}}
\]
This is the intended endpoint of \emph{The Space of Geometry}. The subsequent mathematics begins from a local geometry that is already equipped with displacement, norm, angle, orthogonality, a minimal finite frame, and an explicit symmetry structure.

\appendix
\section{Pauli identities used in the construction}

The Pauli matrices satisfy
\[
\sigma_i\sigma_j=\delta_{ij}I+i\epsilon_{ijk}\sigma_k,
\qquad
\Tr(\sigma_i)=0,
\qquad
\Tr(\sigma_i\sigma_j)=2\delta_{ij}.
\]
For traceless Hermitian operators $A=\bm{a}\cdot\sigmavec$ and $B=\bm{b}\cdot\sigmavec$,
\[
\frac12\Tr(AB)=\bm{a}\cdot\bm{b},
\]
so the Hilbert--Schmidt form restricted to $\Herm_0(2)$ is exactly the standard Euclidean dot product on the Pauli coefficient vector.

\section{Exact regular tetrahedral realization}

An integer-coordinate realization useful for deterministic validation is
\[
\vvec_1=(1,1,1),\quad
\vvec_2=(1,-1,-1),\quad
\vvec_3=(-1,1,-1),\quad
\vvec_4=(-1,-1,1).
\]
It obeys
\[
\sum_{a=1}^{4}\vvec_a=0,
\qquad
\vvec_a\cdot\vvec_a=3,
\qquad
\vvec_a\cdot\vvec_b=-1\quad(a\neq b).
\]
After normalization $\nvec_a=\vvec_a/\sqrt3$,
\[
\nvec_a\cdot\nvec_b=-\frac13\quad(a\neq b),
\]
and
\[
\sum_{a=1}^{4}\nvec_a\nvec_a^T=\frac43I_3.
\]
These exact identities are used by the deterministic repository validators for the finite-cell and informational-convergence branches.

\begin{thebibliography}{9}

\bibitem{LipaTIR}
A. Lipa,
\emph{The Fundamental Theory of Informational Relations},
research repository, 2026.
\newblock \href{https://github.com/AdrianLipa90/The-Fundamental-Theory-of-Informational-Relations}{GitHub repository}.

\bibitem{NielsenChuang}
M. A. Nielsen and I. L. Chuang,
\emph{Quantum Computation and Quantum Information},
10th Anniversary ed., Cambridge University Press, Cambridge, 2010.

\bibitem{RenesEtAl}
J. M. Renes, R. Blume-Kohout, A. J. Scott, and C. M. Caves,
``Symmetric informationally complete quantum measurements,''
\emph{Journal of Mathematical Physics} \textbf{45} (2004), 2171--2180.
\newblock doi: \href{https://doi.org/10.1063/1.1737053}{10.1063/1.1737053}.

\bibitem{Coxeter}
H. S. M. Coxeter,
\emph{Regular Polytopes}, 3rd ed.,
Dover Publications, New York, 1973.

\bibitem{Bloch}
F. Bloch,
``Nuclear induction,''
\emph{Physical Review} \textbf{70} (1946), 460--474.

\end{thebibliography}

\end{document}
